\documentclass[
	%a4paper, % Décommentez pour un format A4 (le format par défaut est US letter)
	11pt, % Taille de police par défaut, peut être 10pt, 11pt ou 12pt
]{resume} % Utilise la classe resume

\usepackage{helvet} % Utilise la police EB Garamond
\usepackage{hyperref}
%------------------------------------------------

\name{Christophe Rouleau-Desrochers} 

\address{
  2424 Main Mall, Vancouver (BC) \\
  (514)·431·4835 \\
  \href{mailto:christophe.rouleaudesrochers@ubc.ca}{christophe.rouleaudesrochers@ubc.ca}
}

\begin{document}

%----------------------------------------------------------------------------------------
%	ÉDUCATION
%----------------------------------------------------------------------------------------

\begin{rSection}{Éducation}
	\textbf{University of British Colombia} \hfill 2024 - Présent\\
	M.Sc. en foresterie
	\smallskip \\
	\textbf{Université du Québec à Montréal} \hfill 2021-2024 \\ 
	B.Sc. en biologie\\ 
	Moyenne générale : 4.30/4.30 (avec mention d'honneur)
	\smallskip \\
	\textbf{Université du Québec à Montréal} \hfill 2020-2021 \\ 
	Certificat en écologie\\ 
	Moyenne générale : 4.19/4.30
\end{rSection}

%----------------------------------------------------------------------------------------
%	Bourses
%----------------------------------------------------------------------------------------

\begin{rSection}{Bourses et distinctions}	
	\textbf{CRSNG} — Bourse d’études supérieures Maîtrise \textbf{(27 000 \$)} \hfill 2025 \\
	\vspace{-2ex}\\
	\textbf{Fondation de l'UQÀM} — Bourse Cocodet en sciences \textbf{(21 375 \$)} \hfill 2023–2024 \\
	\vspace{-2ex}\\
	\textbf{CRSNG} — Bourses de recherche de 1\textsuperscript{er} cycle \textbf{(6 000 \$)} \hfill 2022–2023–2024 \\ 
	\vspace{-2ex}\\
	\textbf{FRQNT} — Supplément de la Bourse de recherche de 1\textsuperscript{er}cycle \textbf{(1 500 \$)} \hfill 2022–2023–2024\\ 
	\vspace{-2ex}\\
	\textbf{CRSNG} — Supplément à l’appui de la diversité de l’effectif dans le secteur forestier \textbf{(5 000 \$)} \hfill 2023 \\ 
	\vspace{-2ex}\\
	\textbf{Faculté des sciences de l’UQÀM} — Liste d’excellence du doyen  \hfill 2023 \\ 
	\vspace{-2ex}\\
	\textbf{Fondation de l'UQÀM} — Bourse Jeanne-D’Arc et Bertin-Trottier en sciences \textbf{(1 500 \$)} \hfill 2023 \\ 
	\vspace{-2ex}\\
	\textbf{Fondation de l'UQÀM} — Bourse du Syndicat des professeurs de l’UQÀM \textbf{(1 500 \$)} \hfill 2023 \\ 
	\vspace{-2ex}\\
	\textbf{Fondation de l'UQÀM} — Bourse de la Fondation Familiale Trottier \textbf{(1 500 \$)} \hfill 2022 \\ 
	\vspace{-2ex}\\
	\textbf{Fondation de l'UQÀM} — Bourse de la Banque TD \textbf{(2 000 \$)} \hfill 2022 
\end{rSection}

%----------------------------------------------------------------------------------------
%	Publications
%----------------------------------------------------------------------------------------
\begin{rSection}{Publications}
Wolkovich, E.M., \textbf{Rouleau-Desrochers, C.}, Garcia de Cortazar-Atauri, I., Walker, M. A., Lacombe, T. (2025). Building a more predictive model of Terroir for the Anthropocene. \textit{Plants, People, Planet} (In press)
	\vspace{1.5ex}\\
\textbf{Rouleau-Desrochers, C.}, Patry, É., Gingras, G., Ucler, J., Vézina-Doré, M., Paillé, O., Moghaddam, Y. Y. (2024) Un été parti en fumée. \textit{Le Point Bio} (non révisé par les pairs)
\end{rSection}

%----------------------------------------------------------------------------------------
%	Expérience de recherche et d’enseignement
%----------------------------------------------------------------------------------------

\begin{rSection}{Enseignement et expérience de recherche}

	\begin{rSubsection}{UBC Faculty of Forestry}{2024-Present}{Graduate Teacher Assistant}{University of British Columbia, Canada}
		\item  FRST-304 Forest stewardship in a changing climate taught by Suzanne Simard for 1217 students (2025)
		\item  FRST-303 Principles of Forest Sciences taught by Sally Aitken for 1962 students (2024) 
	\end{rSubsection}
%------------------------------------------------
	\begin{rSubsection}{Temporal Ecology Lab}{2023-2024}{USRA Undergraduate Researcher}{University of British Columbia, Canada}
		\item Ran a research project on phenological and growth responses of trees to longer growing seasons (2024) 
		\item Involved in a research project on the phenological responses of different tree species drought and defoliation events (2023)
	\end{rSubsection}
%------------------------------------------------
	\begin{rSubsection}{Laboratoire de Pierre Drapeau}{2021-2024}{Assistant de recherche au 1er cycle}{Université du Québec à Montréal, Canada}
		\item Projet sur la phénologie de nidification du Grand pic (2022-2023)
		\item Étude du domaine vital et de l’occupation de l’habitat du Grand pic en forêt boréale mixte (2022)
		\item Participation à un projet sur l’utilisation des poteaux électriques par les pics (2021)
	\end{rSubsection}
%------------------------------------------------ 
	\begin{rSubsection}{Ecological Cascades Lab}{2022}{Assistant de recherche (à distance)}{University of Queensland, Australie}
		\item Participation à un projet de recherche dans les forêts humides des tropiques australiens sur les interactions entre espèces
	\end{rSubsection}
\end{rSection}

%----------------------------------------------------------------------------------------
%	Bénévolat et tutorat
%----------------------------------------------------------------------------------------

\begin{rSection}{Bénévolat et tutorat}
	\begin{rSubsection}{BUDR}{2025}{Mentor}{University of British Columbia, Canada}
		\item Accompagnement de deux étudiant·e·s de premier cycle dans la poursuite de leurs études à la maitrise
	\end{rSubsection}
	\begin{rSubsection}{Faculté des sciences de l’UQAM}{2023-2024}{Tuteur}{Université du Québec à Montréal, Canada}
		\item Soutien à sept étudiant·e·s de première année dans leur cours de biochimie et de statistique
	\end{rSubsection}
	\begin{rSubsection}{Programme Allô}{2022-2024}{Bénévole}{Université du Québec à Montréal, Canada}
		\item Accompagnement de douze étudiant·e·s internationaux sur le campus et transmission de ressources académiques utiles
	\end{rSubsection}
\end{rSection}

%----------------------------------------------------------------------------------------
%	SELECTED OUTREACH
%----------------------------------------------------------------------------------------
\begin{rSection}{VULGARISATION SCIENTIFIQUE}
  \textbf{Tree Spotters — Groupe de science citoyenne} — \textit{Arboretum Arnold, Harvard University} \hfill 2025 \\
  Présentation publique sur l’intérêt des anneaux de croissance pour mieux comprendre le lien entre la phénologie et la dynamique de croissance des arbres. (Boston, MA, États-Unis) 
  
  \textbf{Portes ouvertes de la SCEE} — \textit{Société canadienne d’écologie et d’évolution} \hfill 2024 \\
  Organisation d’une visite de serre et présentation sur l’usage de micro-dendromètres magnétiques pour étudier la croissance des arbres. (Vancouver, BC, Canada)
 
\end{rSection}

%----------------------------------------------------------------------------------------
%	PHOTOGRAPHIE SCIENTIFIQUE
%----------------------------------------------------------------------------------------
\begin{rSection}{PHOTOGRAPHIE SCIENTIFIQUE}
  \textbf{Photo de couverture} – \textit{Proceedings of the National Academy of Sciences} \hfill 2025 \\
  Image présentée en couverture de la section Écologie de l’article \textit{Summer solstice optimizes the thermal growing season}.\\
  \href{https://www.pnas.org/doi/10.1073/iti2325122#sec-4}{https://www.pnas.org/doi/10.1073/iti2325122} 

  \textbf{Concours photo} – \textit{Beaty Biodiversity Museum} \hfill 2024 \\
  Photographie d’un Échasse d’Amérique (South Padre Island, Texas), sélectionnée pour le concours annuel de photos de recherche. \\
  \href{https://explore.beatymuseum.ubc.ca/BRC-photo-competition/25/entries/Christophe-Rouleu-Desrochers/}{https://explore.beatymuseum.ubc.ca/BRC-photo-competition/}

  \textbf{Photo de couverture} – \textit{Le Point Bio} \hfill 2024 \\
  Photographie d’une Aigrette neigeuse choisie pour la couverture du numéro 2024 de \textit{Le Point Bio}.\\
  \href{https://lepointbiologique.com/wp-content/uploads/2024/04/le_point_bio_2024_numerique.pdf}{https://lepointbiologique.com/}\\
\end{rSection}

%----------------------------------------------------------------------------------------
%	Expérience professionnelle
%----------------------------------------------------------------------------------------

\begin{rSection}{Expérience professionnelle}
	\textbf{Fromm Winery (Blenheim, Nouvelle-Zélande)} \hfill 2019-2020 \\ 
	\textit{Viticulteur} \\ 
	\textbf{Restaurant Pastel (Montréal, Canada)} \hfill 2018-2019 \\
	\textit{Assistant-sommelier} \\
	\textbf{Restaurant Toqué! (Montréal, Canada)} \hfill 2017 \\
	\textit{Stagiaire en gestion de cuisine} \\
	\textbf{Winvian Farm (Morris (CT), États-Unis)} \hfill 2017 \\
	\textit{Stagiaire en cuisine}\\
	\textbf{Laurie Raphaël (Montréal, Canada)} \hfill 2016-2017 \\
	\textit{Garde-manger}
\end {rSection}

%----------------------------------------------------------------------------------------
%	Compétences techniques
%----------------------------------------------------------------------------------------

\begin{rSection}{Compétences techniques}
	\begin{tabular}{@{} >{\bfseries}l @{\hspace{6ex}} l @{}}
		Langages informatiques &  R, LaTeX, Git, Bash, Stan, Maxima\\
		Logiciels & Github, Qgis, R, Lightroom, Illustrator
	\end{tabular}
\end{rSection}

\end{document}
